%-------------------------------------------------------------------------------
% layout.tex — LAYOUT & STYLE ONLY (no content)
% Defines document class, page geometry, fonts, colors and the macros used by
% content.tex. Edit this file to change how the CV looks.
% Do not put CV content here. Entry point is resume.tex.
%-------------------------------------------------------------------------------
\documentclass[8pt]{extarticle} % base font size: change 8pt here

\usepackage{geometry} % page size/margins are configured from settings below
\usepackage{fontspec}
\usepackage{paracol}
\usepackage{xcolor}
\usepackage{graphicx}
\usepackage{fontawesome5}
\usepackage{enumitem}
\usepackage{ragged2e}
\usepackage[hidelinks]{hyperref}

\graphicspath{{assets/images/}}

%=============================================================================
%  CUSTOMIZE HERE — every tunable setting lives in this one block.
%  Everything below it only references these constants, so you should not need
%  to edit hex values, font files, sizes, ratios, or spacing anywhere else.
%=============================================================================

%--- Page: paper dimensions and margins --------------------------------------
\newcommand{\cvpaperwidth}{559pt}
\newcommand{\cvpaperheight}{790pt}
\newcommand{\cvmargintop}{22pt}
\newcommand{\cvmarginbottom}{20pt}
\newcommand{\cvmarginleft}{26pt}
\newcommand{\cvmarginright}{19pt}

%--- Fonts: point these at any local .otf/.ttf files under \fontdir ----------
\newcommand{\fontdir}{assets/fonts/}
\newcommand{\fontregular}{SF-Pro-Text-Regular.otf}
\newcommand{\fontbold}{SF-Pro-Text-Semibold.otf}
\newcommand{\fontitalic}{SF-Pro-Text-RegularItalic.otf}
\newcommand{\fontbolditalic}{SF-Pro-Text-SemiboldItalic.otf}
\newcommand{\fontlight}{SF-Pro-Text-Light.otf}

%--- Colors (hex, without the leading #) -------------------------------------
\definecolor{accent}{HTML}{22A39F} % section headers, role line, links, badges
\definecolor{primary}{HTML}{1A1A1A}    % primary text (name, titles, headline, body)
\definecolor{secondary}{HTML}{666666}  % secondary text (education, footer, entry meta)

%--- Font sizes as {size}{leading} in pt -------------------------------------
\newcommand{\sizename}{\fontsize{23pt}{25pt}\selectfont}     % header name
\newcommand{\sizerole}{\fontsize{9pt}{11pt}\selectfont}      % header role line
\newcommand{\sizeheadline}{\fontsize{10pt}{12pt}\selectfont} % header tagline
\newcommand{\sizelabel}{\fontsize{8pt}{8pt}\selectfont}      % section + block labels
\newcommand{\sizebody}{\fontsize{8pt}{10pt}\selectfont}      % normal body text
\newcommand{\sizefooter}{\fontsize{6pt}{8pt}\selectfont}     % small-print footer

%--- Text: line spacing and the name's letter-spacing (tracking) -------------
\newcommand{\cvlinespread}{1.02}
\newcommand{\cvnametracking}{4.0}

%--- Layout: column ratios, gutter, page-2 offset, work split, icon/badge ----
\newcommand{\headerratio}{0.74}  % header row: left column share (0-1)
\newcommand{\bodyratio}{0.505}   % body rows: left column share (0-1)
\newlength{\bodycolsep}\setlength{\bodycolsep}{8pt}         % gutter between body columns
\newlength{\pagetwooffset}\setlength{\pagetwooffset}{68pt}  % top offset aligning page 2 to page 1
\newcommand{\cvworksplit}{6}     % first N jobs go in the left column, rest in the right
\newcommand{\cvsectioniconsize}{9pt}                        % section-header icon height
\newcommand{\cvsectioniconraise}{-0.18em}                   % icon vertical nudge onto the label baseline
\newlength{\cvsectioniconsep}\setlength{\cvsectioniconsep}{6pt}   % gap icon -> section label
\newlength{\cvcoursebadgewidth}\setlength{\cvcoursebadgewidth}{13pt} % course grade badge width

%--- Vertical spacing --------------------------------------------------------
% content.tex never sets \vspace itself; these lengths control every gap and the
% structural macros below insert them automatically. Keep the BETWEEN-block gaps
% larger than the WITHIN-element gaps so sections stay visually separated from
% the lines inside a single element.

% Between blocks:
\newlength{\sectiongap}\setlength{\sectiongap}{6pt}              % between sections (e.g. Introduction -> Experience)
\newlength{\sectionlabelgap}\setlength{\sectionlabelgap}{4pt}    % section label -> its content
\newlength{\headerblockgap}\setlength{\headerblockgap}{6pt}      % between header right-column blocks
\newlength{\headerblocklabelgap}\setlength{\headerblocklabelgap}{2pt} % header block label -> its lines

% Within the header:
\newlength{\namegap}\setlength{\namegap}{6pt}                    % name -> role line
\newlength{\rolegap}\setlength{\rolegap}{2pt}                    % role line -> headline
\newlength{\headerbottomgap}\setlength{\headerbottomgap}{4pt}    % below the whole header
\newlength{\contactlinegap}\setlength{\contactlinegap}{1pt}      % between contact lines

% Within an element:
\newlength{\bodygap}\setlength{\bodygap}{3pt}                    % between consecutive body paragraphs / entries
\newlength{\newlinegap}\setlength{\newlinegap}{3pt}              % between \\ new lines in a body
\newlength{\entrytopgap}\setlength{\entrytopgap}{3pt}            % above an entry title
\newlength{\entrymetagap}\setlength{\entrymetagap}{1pt}          % below an entry meta line
\newlength{\edugap}\setlength{\edugap}{2pt}                      % above an education row
\newlength{\coursegap}\setlength{\coursegap}{1.5pt}             % between course rows
\newlength{\skillgap}\setlength{\skillgap}{2pt}                 % above a skill row
\newlength{\langtoppad}\setlength{\langtoppad}{2pt}            % lead-in above the language pair
\newlength{\langlinegap}\setlength{\langlinegap}{1pt}         % language name -> level line
\newcommand{\cvlangcellwidth}{0.5}                             % language cell width as a share of the column

% Bullet lists (itemize) inside any body:
\newcommand{\cvlistindent}{1em}  % bullet indentation
\newlength{\cvlistitemsep}\setlength{\cvlistitemsep}{0pt}     % between list items
\newlength{\cvlisttopsep}\setlength{\cvlisttopsep}{1pt}       % above/below the whole list
\newlength{\cvlistparsep}\setlength{\cvlistparsep}{0pt}       % between paragraphs in an item
%=============================================================================

%----------------- APPLIES THE SETTINGS ABOVE (do not edit below) -----------------
\geometry{
  paperwidth=\cvpaperwidth, paperheight=\cvpaperheight,
  top=\cvmargintop, bottom=\cvmarginbottom,
  left=\cvmarginleft, right=\cvmarginright
}

\setmainfont{\fontregular}[
  Path            = \fontdir,
  BoldFont        = \fontbold,
  ItalicFont      = \fontitalic,
  BoldItalicFont  = \fontbolditalic
]
\newfontfamily\sflight{\fontlight}[Path=\fontdir]

\color{primary}
\pagestyle{empty}
\setlength{\parindent}{0pt}
\linespread{\cvlinespread}
\renewcommand\normalsize{\sizebody}
\normalsize

% Tracks whether the next \body is the first of its group (no leading gap) or a
% follow-on paragraph (gets \bodygap). Reset by every macro that opens a group.
\newif\ifcvfirstbody
\newcommand{\resetbodygap}{\global\cvfirstbodytrue}
\resetbodygap

%----------------- CONTENT BUCKETS (content.tex fills these) -----------------
% content.tex declares only elements. Each element macro appends its rendered
% form to one of these lists; \makecv (invoked once, from main.tex) lays the
% lists out into sections, columns and pages. No structure lives in content.tex.
\makeatletter
\newcommand{\appendto}[2]{\g@addto@macro#1{#2}}
\makeatother
% Work experience is split across the page-1 column boundary: the first
% \cvworksplit jobs fill the left column, the rest continue in the right column
% (reproducing the original hand-placed layout). \cvworksplit is set above.
\newcounter{cvjobno}
\newcommand{\cvintro}{}
\newcommand{\cvworkleft}{}
\newcommand{\cvworkright}{}
\newcommand{\cvprojects}{}
\newcommand{\cvpubs}{}
\newcommand{\cvedu}{}
\newcommand{\cvcourses}{}
\newcommand{\cvhonours}{}
\newcommand{\cvskills}{}
\newcommand{\cvlangs}{}
\newcommand{\cvsectionfootnote}{}
\newcommand{\cvfootertext}{}

%----------------- STRUCTURAL MACROS -----------------
% Header block: name, role, tagline (left) + location/contacts (right)
% -- Declarative header fields ---------------------------------------------
% content.tex only DECLARES facts with these setters; it never touches the
% document structure. \makecvheader renders them into the two-column header.
%
%   \name{GIACOMO MARCIANI}
%   \role{SENIOR SOFTWARE ENGINEER}   % rendered bold
%   \company{Amazon Web Services}     % rendered normal weight, after "@"
%   \headline{Building High Performance Computing in the Cloud}
%   \location[https://maps.google.com/...]{Boston, MA, United States} % URL optional
%   \email{giacomo.marciani@gmail.com}
%   \linkedin{gmarciani}          % handle only; label/URL are built here
%   \github{gmarciani}            % handle only
%   \makecvheader
%
% Any field left unset is simply omitted from the rendered header.
\newcommand{\cvname}{}     \newcommand{\name}[1]{\renewcommand{\cvname}{#1}}
\newcommand{\cvrole}{}     \newcommand{\role}[1]{\renewcommand{\cvrole}{#1}}
\newcommand{\cvcompany}{}  \newcommand{\company}[1]{\renewcommand{\cvcompany}{#1}}
\newcommand{\cvheadline}{} \newcommand{\headline}[1]{\renewcommand{\cvheadline}{#1}}
% Location accepts an optional URL: \location[https://maps...]{Boston, MA, ...}
% When given, the location text becomes a clickable link.
\newcommand{\cvlocation}{}
\newcommand{\cvlocationurl}{}
\newcommand{\location}[2][]{%
  \renewcommand{\cvlocation}{#2}\renewcommand{\cvlocationurl}{#1}}
\newcommand{\cvemail}{}    \newcommand{\email}[1]{\renewcommand{\cvemail}{#1}}
\newcommand{\cvlinkedin}{} \newcommand{\linkedin}[1]{\renewcommand{\cvlinkedin}{#1}}
\newcommand{\cvgithub}{}   \newcommand{\github}[1]{\renewcommand{\cvgithub}{#1}}

% \iffield{\cvxxx}{code} runs code only when field xxx was set (non-empty).
\newcommand{\iffield}[2]{\ifx\relax#1\relax\else\ifx#1\empty\else#2\fi\fi}

% Render the header from the declared fields.
\newcommand{\makecvheader}{%
  \cvcheckcontact
  \columnratio{\headerratio}%
  \begin{paracol}{2}
  {\bfseries\sizename\color{primary}\addfontfeature{LetterSpace=\cvnametracking}\cvname}\\[\namegap]
  {\sizerole\color{accent}{\bfseries\cvrole}\iffield{\cvcompany}{ @ \cvcompany}}\\[\rolegap]
  {\sizeheadline\color{primary}\cvheadline}
  \switchcolumn
  \cvfirstheaderblocktrue
  \iffield{\cvlocation}{\headerblock{LOCATION}{%
    \ifx\cvlocationurl\empty\cvlocation\else\href{\cvlocationurl}{\cvlocation}\fi}}%
  \ifcvhascontact\headerblock{CONTACTS}{\cvcontactlines}\fi
  \end{paracol}
  \vspace{\headerbottomgap}%
}

% True when at least one contact handle was declared.
\newif\ifcvhascontact
\newcommand{\cvcheckcontact}{%
  \iffield{\cvemail}{\global\cvhascontacttrue}%
  \iffield{\cvlinkedin}{\global\cvhascontacttrue}%
  \iffield{\cvgithub}{\global\cvhascontacttrue}}

% Assemble the contact lines from whichever handles were declared.
\newcommand{\cvcontactlines}{%
  \iffield{\cvemail}{\href{mailto:\cvemail}{\cvemail}\\[\contactlinegap]}%
  \iffield{\cvlinkedin}{\href{https://linkedin.com/in/\cvlinkedin}{\cvlinkedin\ (LinkedIn)}\\[\contactlinegap]}%
  \iffield{\cvgithub}{\href{https://github.com/\cvgithub}{\cvgithub\ (GitHub)}}%
}

% A labelled block for the header right column (e.g. Location, Contacts).
% #1 = label  #2 = lines. A leading gap separates stacked blocks; the first
% block in the column gets none.
\newif\ifcvfirstheaderblock
\newcommand{\headerblock}[2]{%
  \ifcvfirstheaderblock\global\cvfirstheaderblockfalse\else\vspace{\headerblockgap}\fi
  {\color{accent}\sizelabel #1}\\[\headerblocklabelgap]#2\par}

%----------------- STRUCTURAL SKELETON (used only by \makecv) -----------------
% content.tex never calls any of these; \makecv assembles the document.
% Open / close the two independent body columns
\newcommand{\bodystart}{%
  \columnratio{\bodyratio}\setlength{\columnsep}{\bodycolsep}\begin{paracol}{2}}
\newcommand{\bodyend}{\end{paracol}}
% Switch from the left body column to the right one
\newcommand{\rightcolumn}{\switchcolumn}
% Page break that also applies the page-2 top alignment offset
\newcommand{\cvnewpage}{\newpage\vspace*{\pagetwooffset}}

% Section header: teal icon (from assets/images/) + teal uppercase label
% #1 = icon file (in assets/images/), #2 = label
\newcommand{\cvsection}[2]{%
  \resetbodygap
  \vspace{\sectiongap}%
  \noindent\raisebox{\cvsectioniconraise}{\includegraphics[height=\cvsectioniconsize]{#1}}\hspace{\cvsectioniconsep}%
  {\color{accent}\sizelabel #2}\par
  \vspace{\sectionlabelgap}%
}

%----------------- ELEMENT RENDERERS (private) -----------------
% These typeset one element. The public element macros below only APPEND a
% call to one of these into a section bucket; nothing is typeset until \makecv.

% Entry: bold title, secondary meta line. Optional URL makes the title a clickable
% teal link.  \entry[url]{Title}{Meta}
\newcommand{\entry}[3][]{%
  \resetbodygap
  \vspace{\entrytopgap}%
  {\bfseries\ifx\relax#1\relax #2\else\href{#1}{\color{accent}#2}\fi}\par
  {\color{secondary}#3}\par
  \vspace{\entrymetagap}%
}
% Entry variant where linked titles use primary color instead of accent.
\newcommand{\entryprimary}[3][]{%
  \resetbodygap
  \vspace{\entrytopgap}%
  {\bfseries\ifx\relax#1\relax #2\else\href{#1}{\color{primary}#2}\fi}\par
  {\color{secondary}#3}\par
  \vspace{\entrymetagap}%
}
% A \\ inside any element body starts a new line separated by \newlinegap.
\newcommand{\bodynewline}{\par\vspace{\newlinegap}}
% Justified body paragraph. Consecutive paragraphs are separated by \bodygap;
% the first after a section/entry gets no leading gap. Within the text, \\
% inserts a \newlinegap-spaced new line.
\long\def\body#1{%
  \ifcvfirstbody\global\cvfirstbodyfalse\else\vspace{\bodygap}\fi
  {\justifying\let\\\bodynewline #1\par}}
% Education row: bold degree, muted school, muted years+grade.
\newcommand{\renderedu}[3]{%
  \resetbodygap
  \vspace{\edugap}{\bfseries #1}\par
  {\color{secondary}#2}\par
  {\color{secondary}#3}\par}
% Course row: name (left) + teal grade badge (right).
\newcommand{\rendercourse}[2]{%
  \noindent#1\hfill\colorbox{accent}{\makebox[\cvcoursebadgewidth][c]{\color{white}\bfseries#2}}\par
  \vspace{\coursegap}}
% Skill row: bold category + description.
\newcommand{\renderskill}[2]{%
  \vspace{\skillgap}\noindent{\bfseries #1}\quad #2\par}
% One language proficiency as a half-column cell; consecutive cells sit side by
% side, so declaring two languages reproduces the original paired row.
\newcommand{\renderlanguage}[2]{%
  \begin{minipage}[t]{\cvlangcellwidth\columnwidth}{\bfseries #1}\\[\langlinegap]{\color{secondary}#2}\end{minipage}}
% Small-print footer.
\newcommand{\cvfooter}[1]{%
  \vfill{\sflight\sizefooter\color{secondary} #1\par}}

%----------------- ELEMENT DECLARATIONS (content.tex uses only these) --------
% Each macro appends its element to the relevant section bucket. content.tex is
% a flat list of these calls: no sections, columns, pages, or body plumbing.
% In any body/description, \\ starts a new line spaced by \newlinegap.
%
%   \introduction{First line.\\ Second line.\\ Third line.}
%   \job[url]{Title}{Employer}{Years}{Description with a\\ new line}
%   \project[url]{Title}{Meta}{Description with a\\ new line}
%   \sectionfootnote{Trailing italic note under a section.}
%   \publication[url]{Title}{Venue}     % URL optional; title becomes a link
%   \education{Degree}{School}{Years and grade}
%   \course{Name}{Grade}
%   \skill{Category}{Comma-separated items}
%   \lang{Name}{Level}
%   \footer{Small-print footer text.}
% Introduction: a single element. Like every body, \\ starts a new line
% (spaced by \newlinegap). content.tex declares exactly one \introduction.
\long\def\introduction#1{\appendto\cvintro{\body{#1}}}
% Jobs route to the left column bucket until \cvworksplit is reached, then to
% the right column bucket — so content.tex just lists jobs in order.
\newcommand{\job}[5][]{%
  \stepcounter{cvjobno}%
  \ifnum\value{cvjobno}>\cvworksplit\relax
    \appendto\cvworkright{\entry[#1]{#2}{#3, #4}\body{#5}}%
  \else
    \appendto\cvworkleft{\entry[#1]{#2}{#3, #4}\body{#5}}%
  \fi}
\newcommand{\project}[4][]{\appendto\cvprojects{\entryprimary[#1]{#2}{#3}\body{#4}}}
\newcommand{\sectionfootnote}[1]{\appendto\cvsectionfootnote{\body{\color{secondary}\itshape #1}}}
\newcommand{\publication}[3][]{\appendto\cvpubs{\entryprimary[#1]{#2}{#3}}}
\newcommand{\honour}[2]{\appendto\cvhonours{\entry{#1}{#2}}}
\newcommand{\education}[3]{\appendto\cvedu{\renderedu{#1}{#2}{#3}}}
\newcommand{\course}[2]{\appendto\cvcourses{\rendercourse{#1}{#2}}}
\newcommand{\skill}[2]{\appendto\cvskills{\renderskill{#1}{#2}}}
\newcommand{\lang}[2]{\appendto\cvlangs{\renderlanguage{#1}{#2}}}
\newcommand{\footer}[1]{\renewcommand{\cvfootertext}{#1}}

%----------------- DOCUMENT ASSEMBLY -----------------
% Renders the whole CV from the declared header fields and section buckets.
% Fixed placement: page 1 (Intro+Work | Projects+Publications),
% page 2 (Education+Courses | Honours+Skills+Languages).
\newcommand{\makecv}{%
  \makecvheader
  % --- Page 1 ---
  \bodystart
    \cvsection{introduction}{INTRODUCTION}\cvintro
    \cvsection{work}{WORK EXPERIENCE}\cvworkleft
  \rightcolumn
    \cvworkright
    \cvsection{projects}{MAIN PROJECTS}\cvprojects\cvsectionfootnote
    \cvsection{publications}{PUBLICATIONS}\cvpubs
  \bodyend
  % --- Page 2 ---
  \cvnewpage
  \bodystart
    \cvsection{education}{EDUCATION}\cvedu
    \cvsection{courses}{MAIN COURSES}\cvcourses
  \rightcolumn
    \cvsection{honours}{HONOURS}\cvhonours
    \cvsection{skills}{MAIN SKILLS}\cvskills
    \cvsection{languages}{LANGUAGES}\vspace{\langtoppad}\noindent\cvlangs\par
  \bodyend
  \cvfooter{\cvfootertext}%
}

\setlist[itemize]{leftmargin=\cvlistindent, itemsep=\cvlistitemsep, topsep=\cvlisttopsep, parsep=\cvlistparsep}
